\documentclass[12pt]{article}
\usepackage{amsmath, amssymb}
\usepackage{geometry}
\usepackage{booktabs}
\usepackage{listings}
\usepackage{xcolor}

\geometry{margin=1in}

\title{Problem Set 4}
\author{Qi Liu (25340516)}
\date{}

% ---------- R code style (matches PS3 answer key) ----------
\definecolor{codegreen}{rgb}{0,0.6,0}
\definecolor{codered}{rgb}{0.6,0,0}
\definecolor{codepurple}{rgb}{0.58,0,0.82}
\definecolor{backcolour}{rgb}{0.95,0.95,0.95}

\lstdefinestyle{RStyle}{
	language=R,
	backgroundcolor=\color{backcolour},
	commentstyle=\color{codegreen},
	keywordstyle=\color{codepurple},
	stringstyle=\color{codered},
	basicstyle=\ttfamily\small,
	breaklines=true,
	numbers=left,
	numberstyle=\tiny\color{gray},
	frame=single,
	tabsize=2,
	showstringspaces=false
}

\begin{document}
	
	\maketitle
	
	\section*{Question 1}
	
	We are interested in how income and occupational type are related to
	occupational prestige. We focus on whether professional occupations have
	higher prestige and whether the effect of income differs between
	professional and non-professional jobs.
	
	\subsection*{1. Create the professional dummy variable}
	
	We create a dummy variable that equals 1 for professional occupations and
	0 for blue- and white-collar occupations.
	
	\begin{lstlisting}[style=RStyle]
		# Load data
		library(car)
		data(Prestige)
		
		# Create dummy variable: professional = 1 if type == "prof", else 0
		Prestige$professional <- ifelse(Prestige$type == "prof", 1, 0)
		
		# Check coding
		table(Prestige$professional, Prestige$type)
	\end{lstlisting}
	
	\subsection*{2. Regress prestige on income, professional, and interaction}
	
	We estimate the following model:
	\[
	prestige_i = \beta_0 + \beta_1 income_i + \beta_2 professional_i
	+ \beta_3 (income_i \cdot professional_i) + \varepsilon_i.
	\]
	
	The regression is run in R as follows:
	
	\begin{lstlisting}[style=RStyle]
		# Run regression with interaction
		m1 <- lm(prestige ~ income * professional, data = Prestige)
		summary(m1)
	\end{lstlisting}
	
	Table~\ref{tab:q1} reports the regression results.
	
	\begin{table}[h!]
		\centering
		\caption{Outcome variable is prestige. Predictors are income,
			professional, and their interaction.}
		\label{tab:q1}
		\begin{tabular}{l c}
			\toprule
			& Model 1 \\
			\midrule
			(Intercept)               & 21.1423*** \\
			& (2.8044)   \\
			income                    & 0.0031709*** \\
			& (0.0004993) \\
			professional              & 37.7813*** \\
			& (4.2483)   \\
			income:professional       & -0.0023257*** \\
			& (0.0005675) \\
			\midrule
			R-squared                 & 0.7872     \\
			Adj.\ R-squared           & 0.7804     \\
			Num.\ obs.                & 98         \\
			\bottomrule
		\end{tabular}
		
		\begin{flushleft}
			{\footnotesize Standard errors in parentheses.
				***p$<$0.001, **p$<$0.01, *p$<$0.05.}
		\end{flushleft}
	\end{table}
	
	\subsection*{3. Prediction equation}
	
	Using the estimated coefficients from Table~\ref{tab:q1}, the fitted equation is
	\[
	\widehat{prestige} =
	21.1423
	+ 0.0031709 \cdot income
	+ 37.7813 \cdot professional
	- 0.0023257 (income \cdot professional).
	\]
	
	\subsection*{4. Interpretation of income coefficient}
	
	The coefficient on income is 0.0031709. This represents the marginal
	effect of income on prestige for non-professional occupations
	(professional = 0). Holding professional status fixed at 0, a one-unit
	(one dollar) increase in income is associated with an average increase of
	0.00317 points in prestige. A 1{,}000-unit increase in income
	corresponds to an increase of about 3.17 points in prestige.
	
	\subsection*{5. Interpretation of professional coefficient}
	
	The coefficient on the professional dummy is 37.7813. Because the model
	includes an interaction between income and professional, this coefficient
	captures the difference in predicted prestige between professional and
	non-professional occupations when income equals 0. At income = 0,
	professional occupations are predicted to have about 37.78 more prestige
	points than non-professional occupations.
	
	\subsection*{6. Effect of a 1{,}000-unit income increase for professionals}
	
	With an interaction term, the marginal effect of income is
	\[
	\frac{\partial \widehat{prestige}}{\partial income}
	= \beta_1 + \beta_3 \cdot professional.
	\]
	
	For professional occupations (professional = 1), this becomes
	\[
	ME_{income|prof=1} = 0.0031709 - 0.0023257 = 0.0008452.
	\]
	
	Thus, for professionals, a 1{,}000-unit increase in income changes prestige by
	\[
	1000 \times 0.0008452 \approx 0.8452.
	\]
	
	The corresponding R code is:
	
	\begin{lstlisting}[style=RStyle]
		beta <- coef(m1)
		marg_income_prof1_per1 <- beta["income"] + beta["income:professional"]
		marg_income_prof1_1000 <- marg_income_prof1_per1 * 1000
	\end{lstlisting}
	
	This shows that income has a positive but comparatively small effect on
	prestige among professional occupations.
	
	\subsection*{7. Effect of switching to professional at income = 6000}
	
	The effect of changing from non-professional (professional = 0) to
	professional (professional = 1) at a given income level is
	\[
	\Delta prestige = \beta_2 + \beta_3 \cdot income.
	\]
	
	At income = 6000, we obtain
	\[
	\Delta prestige = 37.7813 + 6000(-0.0023257) = 23.8270.
	\]
	
	In other words, at an income level of 6{,}000, professional occupations are
	predicted to have about 23.83 more prestige points than non-professional
	occupations.
	
	The calculation is implemented in R as:
	
	\begin{lstlisting}[style=RStyle]
		effect_prof_switch_6000 <-
		beta["professional"] + 6000 * beta["income:professional"]
	\end{lstlisting}
	
	%========================================================
	% Question 2
	%========================================================
	\section*{Question 2}
	
	We now consider an experiment on campaign yard signs in a gubernatorial
	election. The outcome variable is the vote share of Cuccinelli, and the
	key predictors are indicators for whether a precinct was assigned yard
	signs and whether it was adjacent to precincts with yard signs. The
	regression model is
	\[
	voteshare_i = \beta_0 + \beta_1 assigned_i + \beta_2 adjacent_i
	+ \varepsilon_i.
	\]
	
	The reported coefficients are summarized in Table~\ref{tab:q2}.
	
	\begin{table}[h!]
		\centering
		\caption{Effect of campaign yard signs on Cuccinelli's vote share.}
		\label{tab:q2}
		\begin{tabular}{l c}
			\toprule
			& Model 1 \\
			\midrule
			(Intercept)                     & 0.302*** \\
			& (0.011)  \\
			Assigned lawn signs            & 0.042**  \\
			& (0.016)  \\
			Adjacent to lawn signs         & 0.042*** \\
			& (0.013)  \\
			\midrule
			R-squared                      & 0.094    \\
			Num.\ obs.                     & 131      \\
			\bottomrule
		\end{tabular}
		
		\begin{flushleft}
			{\footnotesize Standard errors in parentheses.
				***p$<$0.01, **p$<$0.05.}
		\end{flushleft}
	\end{table}
	
	\subsection*{1. Do yard signs affect vote share?}
	
	We first test whether having yard signs in a precinct affects vote
	share. The null and alternative hypotheses for the coefficient on
	assigned are:
	\[
	H_0: \beta_1 = 0, \qquad H_1: \beta_1 \neq 0.
	\]
	
	The estimated coefficient is 0.042 with a standard error of 0.016. The
	t-statistic is
	\[
	t = \frac{0.042}{0.016} \approx 2.63.
	\]
	
	Using a two-sided test with significance level 0.05, the critical value
	is about 1.98. Since the absolute t-statistic is larger than 1.98, we
	reject the null hypothesis and conclude that yard signs have a
	statistically significant effect on vote share. Substantively, precincts
	that were assigned yard signs are predicted to give Cuccinelli about
	4.2 percentage points more of the vote than similar precincts without
	signs.
	
	\subsection*{2. Do adjacent precincts show an effect?}
	
	We next consider the effect of being adjacent to precincts with yard
	signs. The estimated coefficient is again 0.042, with a standard error
	of 0.013. The t-statistic is
	\[
	t = \frac{0.042}{0.013} \approx 3.23.
	\]
	
	This value is also larger than 1.98 in absolute value, so we reject the
	null hypothesis that the coefficient on adjacent is zero. We conclude
	that precincts adjacent to yard signs also show a statistically
	significant increase in Cuccinelli's vote share of about 4.2 percentage
	points.
	
	\subsection*{3. Interpretation of the intercept}
	
	The intercept is estimated as 0.302. This represents the expected vote
	share for Cuccinelli in precincts that neither had yard signs nor were
	adjacent to precincts with yard signs. In the absence of any yard sign
	exposure, Cuccinelli is predicted to receive about 30.2 percent of the
	vote.
	
	\subsection*{4. Model fit and importance of yard signs}
	
	The model has R-squared equal to 0.094, meaning that the two yard sign
	variables explain about 9.4 percent of the variation in vote share
	across precincts. This indicates that while yard signs and adjacency
	have statistically detectable effects, most of the variation in vote
	share is driven by other factors not included in the model, such as
	demographics, partisanship, candidate quality, or other campaign
	activities. Yard signs matter, but they are only one of many
	determinants of electoral outcomes.
	
\end{document}
